% 章节结构可按课题调整，不是学院规定的固定章数。
% 写作建议：先明确具体工程问题，再组织研究现状、问题缺口和研究目标。
\chapter{绪论}
\section{研究背景与工程问题}
本节阐述研究背景、工程需求及研究意义，明确拟解决的问题、研究对象和适用范围。涉及工程规模、水质或运行数据时，应注明来源及时间范围。

污水处理研究涉及水质特征、处理工艺与资源回收等环节\supercite{metcalf2014}。排水系统设计中，工程参数与适用条件应结合相应标准确定\supercite{gb50014}。
% 以上两句展示图书与标准的引用写法，按实际课题替换。
\section{国内外研究现状}
\subsection{相关技术与方法}
本节围绕研究问题梳理代表性文献，比较不同技术与方法的适用条件、主要结论及局限。例如，厌氧氨氧化研究报道了厌氧条件下的铵转化现象\supercite{mulder1995}。
% 该句用于展示期刊论文的引用写法，不限定模板的研究主题。

文献按正文首次出现的顺序编号，引用采用方括号上标，连续引用格式如\supercite{tongji2025guide,tongji2026degree}。
% 上句的学校文件用于展示中文机构作者及网页文献，按正式内容替换。
% 正常引用：研究结论的相应表述\supercite{文献键}。
% 多篇引用：\supercite{key1,key2,key3}；指定页码：\supercite[25--28]{key1}。
% 不手工输入 [1]，不把文献表作为普通文本粘贴。
\subsection{现有研究的不足}
本节概括已有研究尚未解决的问题，说明其与本研究工程对象的联系，明确研究切入点。
\section{研究目标与主要内容}
本节列明研究目标、拟回答的问题及主要研究内容，并说明各项内容之间的逻辑关系。
\section{技术路线与论文结构}
本节说明研究方法、试验验证与工程评价的组织方式，并概述论文各章内容。
